\section{The buffered canonical WAD is a different object}
\label{sec:buffer}
Let $K_0,K_1,K_2$ be the three
\emph{actual} little Julia intervals in
\eqref{eq:tree}.  The minimal
four-boundary KKL surface is
\begin{equation}\label{eq:four}
 Y_{K,p}=D(0,R_p)\setminus(K_0\cup K_1\cup K_2).
\end{equation}
Its four \emph{whole ideal boundary
components}, rather than ten or
twenty individual cut atoms, are
$I_{\rm out},I_{K_0},I_{K_1},I_{K_2}$.
For an A.1 admissible unoriented
simple ideal class $\gamma$,
choose its two \emph{universal-cover}
ideal interval lifts, and put
$q_\gamma$ equal to their full
quadrilateral extremal width.
KKL \cite[Appendix A.2.3--A.2.5]{KKL2026}
defines the buffered weight
$w_\gamma=(q_\gamma-2)_+$ and
$\lVert W_{\rm can}\rVert=\sum_\gamma
w_\gamma$.  The exceptional
unbuffered annulus convention
does not apply to \eqref{eq:four}.
Self-endpoint types can use
distinct lifts and must not be
summarily discarded.

\begin{lemma}[Uniform cyclic-cover exclusion test]
\label{lem:peripheral}
Let $I$ be any minimally
marked nonsimply connected
peripheral ideal circle
of a finite-type hyperbolic
surface which is not the
exceptional marked annulus.
Write $\W(I)$ for the
width of its peripheral
cyclic cover.  For \emph{every}
A.1 admissible class $\gamma$
incident on $I$, including
essential self-endpoint
classes using different lifts,
\begin{equation}\label{eq:peripheral}
 q_\gamma\le
 \frac{K(\sqrt{1-e^{-\pi\W(I)}})}
 {K(e^{-\pi\W(I)/2})}.
\end{equation}
If $\pi\W(I)<\log(17+12\sqrt2)$,
then $q_\gamma<2$ and
$w_\gamma=0$ for all such
classes, with \emph{no}
winding-word cutoff.
\end{lemma}
\begin{proof}
Uniformize the universal
cover as $\mathbb H$ with a
lift of $I$ equal to
$(0,\infty)$ and its
primitive stabilizer
$\tau(z)=\lambda z$,
$\lambda>1$.  Writing
$\ell=\log\lambda$,
the logarithmic cyclic
quotient is a strip
of height $\pi$ and
circumference $\ell$.
Its transverse width is
$\ell/\pi$, not the reciprocal:
$\ell=\pi\W(I)$.

Another ideal lift
is, after normalization,
$J=(-b,-a)$ with $0<a<b$.
The translates $J$ and
$\tau J$ have disjoint
interiors: different
lifts of ideal boundary
components do not overlap.
If a translate stabilized
$J$ itself, its peripheral
stabilizer would coincide
with that of $I$,
forcing the same
boundary lift and an
A.1-boundary-subordinate
class, which is excluded.
Thus $b/a\le\lambda$.
For ideal intervals
$(0,\infty)$ and $(-b,-a)$
in $\mathbb H$, the
Schwarz--Christoffel
four-point formula gives
$k^2=1-a/b$ and
$q_\gamma=K(k)/K(\sqrt{1-k^2})$.
Since this ratio increases
strictly with $k$, we
obtain \eqref{eq:peripheral}.
The descending quadratic
Landen identity, applied
at $k_0=1/\sqrt2$, gives
\[
 \frac{K(\sqrt{12\sqrt2-16})}
      {K(\sqrt{17-12\sqrt2})}=2,
 \qquad
 1-(12\sqrt2-16)=17-12\sqrt2
  =(17+12\sqrt2)^{-1}.
\]
Hence the strict peripheral
length cutoff stated in
the lemma implies
$q_\gamma<2$.
The continuum inputs are
the uniformization,
peripheral stabilizers,
SC quadrilateral formula
and Landen's identity,
not Python theorems.
\end{proof}

\begin{theorem}[All outer-incident buffered weights vanish
for $p\ge11$]\label{thm:buffer}
For $p=11,12$ and
\emph{every} integer $p\ge13$,
every KKL buffered class on
$Y_{K,p}$ incident on
$I_{\rm out}$ has $q_\gamma<2$
and $w_\gamma=0$,
including all possible
essential outer self-classes.
Moreover
\[
 0\le\lVert W_{\rm can}\rVert<619/200.
\]
Only six \emph{inner
endpoint types} remain
possible: $(K_i,K_j)$
for $0\le i\le j\le2$.
This is \emph{not}
an enumeration or exclusion
of all infinitely many
inner isotopy classes.
\end{theorem}
\begin{proof}
In the positive scale
$\xi=-A_pz$ a real segment
$[l,r]$, with midpoint
$m$ and half-length $a$,
has Joukowski chart
$\xi=m+(a/2)(w+w^{-1})$.
An embedded ellipse
$|w|=t>1$ has horizontal
semiaxis $(a/2)(t+t^{-1})$.
If its entire annulus
is contained in the
physical complement
of all four holes,
its essential cyclic
annulus lifts into the
corresponding peripheral
cover.  The explicit
potential $\log|w|/\log t$,
extended by constants,
has Dirichlet energy
$2\pi/\log t$.
The elementary annular
capacity comparison
therefore yields
\begin{equation}\label{eq:collar}
 \W(I)\le\frac{2\pi}{\log t},
 \qquad \pi\W(I)\le
 \frac{2\pi^2}{\log t}.
\end{equation}
For the \emph{outer}
component, use the whole
tree carrier
$[x_{2p},x_p]$ to build
the surrounding ellipse;
the carrier is not
misidentified as a
fourth Julia boundary.
The outward dyadic
finite chart checks give
an embedded outer ellipse
with $t=300$ at $p=11$
and $t=512$ at $p=12$,
both with
$\pi\W(I_{\rm out})<7/2$.
For all $p\ge13$ the
largest normalized
coordinate box occurs at
$p=13$, where the carrier
has $|m|<2$, $0<a<2$ and
$\lambda_pR_p>(3/20)2^p$.
With $t=1024$,
\[
 |m|+\tfrac a2(t+t^{-1})
 <2+1024+1/1024
 <(3/20)2^{13}
 \le\lambda_pR_p.
\]
The full ellipse, not a
sampled boundary point,
thus embeds for every
$p\ge13$, and
\[
 \pi\W(I_{\rm out})\le
 \frac{2\pi^2}{\log 1024}
 =\frac{\pi^2}{5\log2}<3
\]
using $\pi^2<10$
and $\log2>2/3$.
For $p=11,12$ the
simple exact sufficient
test $e^{7/2}<(33/20)^7
<169/5<17+12\sqrt2$
makes Lemma~\ref{lem:peripheral}
apply.  For the tail,
$e^3<27<17+12\sqrt2$
does the same.  All
four outer-incident
unoriented endpoint
types therefore vanish,
regardless of winding.

There are also disjoint
whole ellipses around
the individual $K_0,K_1,K_2$,
with $t$ respectively
$11,49,25$.  The exact
finite and all-$p$
physical-domain checks
include exclusion of
\emph{both other holes}
from each collar.
Equation \eqref{eq:collar}
gives the strict widths
\[
 \W(I_{K_0})<2.621,\quad
 \W(I_{K_1})<1.615,\quad
 \W(I_{K_2})<1.953.
\]
KKL Appendix A.2.3--A.2.5
bounds the total canonical
buffered weight, counting
returning arcs twice in
the local sums, by half
the sum of the peripheral
widths.  Since the outer
weights vanish,
\[
 0\le\lVert W_{\rm can}\rVert
 <\tfrac12(2.621+1.615+1.953)
 =3.0945<619/200.
\]
The remaining six
inner endpoint \emph{types}
can contain infinitely
many separate classes;
their classwise $q_\gamma$
and even the positivity
of the complete WAD
remain open.
\end{proof}

\section{Crest transport, nesting and global QC: precise limits}
\label{sec:limits}
\subsection{Finite collared four-push crest descendants}
The finite crest-transport problem uses
\emph{chosen nonmaximal circular parents},
not $P_p$ or $W_{\rm can}$.
At $p=3$ the parent is the native
Wave centered at $-9/64$,
$3/256<r<1/64$; for $p\ge4$
its chosen center is $-3x_p/4$
and radii $x_p/16<r<(p-1)x_p/(8p)$.
The parent conformal width is
$\log R_p^{\rm chosen}/(2\pi)$,
where $R_3^{\rm chosen}=4/3$
and $R_p^{\rm chosen}=2(p-1)/p$
for $p\ge4$.  The four
successive \emph{proper}
$\iota^*$ restrictions and
separate covering pushes
select two source terminal
families with endpoints
$(S_0,S_3)$ and
$(S_{17},S_0)$; their
target bases are
$(I_0,I_1,I_2,I_3)$
and $(I_7,I_8,I_9,I_0)$,
and their closed crests
$\overline{I_1\cup I_2}$
and $\overline{I_8\cup I_9}$
meet \emph{opposite sides}
$I_1,I_9$ of the \emph{same}
old edge $E_0$.
There is no assertion for
every third-stage path.

Let $\theta_p(r)=\arccos u_p(r)$,
where $u_p(r)\in(-1/8,1/4)$ is
the unique zero of
\[
 \Re G_p\!\left(d_p+r
 (u+i\sqrt{1-u^2})\right)=0.
\]
On a certified terminal
radius interval $(a,b)$
the conformal foliation
calculation gives the
\emph{selected} descendant
width and ratio
\begin{equation}\label{eq:crest}
 \W(S^\pm_{4,p})=
 \int_a^b\frac{dr}{r\,\theta_p(r)},
 \qquad
 \frac{\W(S^\pm_{4,p})}{\W(S_p)}
 =\frac{2\pi}{\log R_p^{\rm chosen}}
       \int_a^b\frac{dr}{r\,\theta_p(r)}.
\end{equation}
This follows from the
conformal logarithmic
strip with variable
angular height:
Cauchy--Schwarz
provides the extremal
density $1/\theta_p(e^x)$.
The archival certificate
bounds \emph{whole}
rational radius bins,
not sampled radii, using
the square-root-free
real/imaginary recurrence
\[
 X_{j+1}=X_j^2-Y^2B_j^2+c_p,\quad
 B_{j+1}=2X_jB_j,\quad
 Y^2=r^2(1-u^2).
\]
The relative cut and
full complex gate checks
make the family a
legitimate selected
KKL proper-path
subfamily.

For the maximal connected
\emph{circular collar only}
among those chosen
parent radii, the strict
finite outward enclosures
from the archive are:
\begin{center}
\begin{tabular}{rcc}
\toprule
$p$& lower bound $c_p^*$& upper bound $c_p^*$\\\midrule
3&0.98017&0.98038\\
4&0.62810&0.62818\\
5&0.52758&0.52764\\
6&0.47562&0.47566\\
8&0.43006&0.43010\\
12&0.39680&0.39684\\
13&0.39229&0.39233\\
16&0.38261&0.38265\\
\bottomrule
\end{tabular}
\end{center}
At $p=3$ the much
\emph{smaller fixed subcollar}
in the first transport
audit gives a \emph{different}
ratio in $(0.5676,0.5684)$;
it must not be conflated
with the full circular
optimum $(0.98017,0.98038)$.
These are rigorous
finite interval statements
about specifically
selected collars, and
not a canonical supremum
over all KKL Waves.

\begin{question}[Crest-transport limit]\label{q:crest}
Numerics suggest that
$c_p^*$ tends to the
positive constant
$c_\infty\approx0.3470049860435258$.
The proposed normalized
quadratic map is
$Q_\kappa(t)=1-\kappa t^2$,
where $\kappa\approx1.7903274919993457$
solves
$Q_\kappa^6(0)+Q_\kappa^3(0)=0$
on the distinguished branch.
The proposed first-order
law is
\[
 c_p^*=
 \frac{A_\infty}{\log(2-2/p)}
       +O(4^{-p}),\quad
 A_\infty\approx0.2405255277163131.
\]
Neither the convergence,
the $O(4^{-p})$ error,
nor a uniform all-period
selected q-side collar
survival is proved by
the finite audits.
The parameter asymptotic
\eqref{eq:tip-rate}
does \emph{not} alone
prove this stronger
four-stage
complex-domain and
prime-end statement.
\end{question}

\subsection{Depth two and the unique certified depth-three row}
For a \emph{realized} primitive
tuning word
$(p_1,\ldots,p_d)$
the depth-$j$ return is a
\emph{central degree-two
polynomial-like restriction}
of $f_c^{n_j}$ with
$n_j=\prod_{i=1}^j p_i$,
not the unrestricted
degree-$2^{n_j}$ polynomial.
For the depth-two words
$(p_1,p_2)\in\{3,4,5\}^2$,
the published archive
certifies nine locally
isolated real parameter
root cells, signed cut
orders, central degree-two
returns, a \emph{separately
normalized} third-return
Rouch\'e gap, and the
following strict \emph{lower
witness} and \emph{full-path
upper envelope}:
\begin{center}
\begin{tabular}{rcc}
\toprule
$(p_1,p_2)$ & $W^{\rm rad}$ lies in&
$Q_{\rm plane}$ lies in\\\midrule
$(3,3)$&(0.25224,0.25225)&(0.76475,0.76476)\\
$(3,4)$&(0.25854,0.25855)&(0.77250,0.77251)\\
$(3,5)$&(0.25999,0.26000)&(0.77430,0.77431)\\
$(4,3)$&(0.25179,0.25180)&(0.76418,0.76419)\\
$(4,4)$&(0.25848,0.25849)&(0.77243,0.77244)\\
$(4,5)$&(0.25999,0.26000)&(0.77431,0.77432)\\
$(5,3)$&(0.25170,0.25171)&(0.76406,0.76407)\\
$(5,4)$&(0.25847,0.25848)&(0.77242,0.77243)\\
$(5,5)$&(0.26000,0.26001)&(0.77431,0.77432)\\
\bottomrule
\end{tabular}
\end{center}
Thus $0.25<W^*<0.775$
holds at these \emph{nine}
physical depth-two charts,
using the same analytic
KKL and prime-end inputs
as at level one.  The first
return is straightened
to the designated suffix
only with a separate
signed kneading rigidity
identification
\cite[Theorem II.10.1]{deMeloVanStrien1993};
equating two independent
tip equations
$P_{p_1}=P_{p_1p_2}=0$
would be incorrect.
The data are finite exact
rational \emph{witness and
envelope} bounds, not
a characterization of
the optimum.

For depth three only the
realized word $(3,3,3)$
has the archived
complete independent
certificate:
$n_j=3,9,27$,
the nested target radii
$1/4,1/40,1/250$,
the three initial
Rouch\'e gaps strictly
$>0.127470,\ 0.006346,\
0.003900$, a \emph{separately
normalized} $81$-return
gap $>0.172673$,
and
\[
 W^{\rm rad}>0.252105,\qquad
 Q_{\rm plane}<0.764704.
\]
Consequently the analogous
physical marked supremum
has $0.25<W^*<0.775$
\emph{for this one
depth-three row}.
None of these exact
finite recurrences proves
an all-word or infinite
tower theorem.

\subsection{Complex bounds, local hybrid maps and the missing global mark}
Lyubich--Yampolsky
\cite[p.~1220, Complex Bounds Theorem]{LyubichYampolsky1997}
gives an absolute positive
\emph{optimal} fundamental
annulus modulus for
every actually realized
finite non-doubling
renormalization of a real
quadratic, subject to its
explicit last-doubling
near-cusp exception.
This is \emph{not}
Theorem 1.1 restricted
to bounded essential period,
and it is not a modulus
bound for the specific
physical round gates
used here.
Douady--Hubbard
\cite[Theorem 1]{DH1985}
straightens each
appropriate quadratic-like
restriction \emph{locally}
around its little filled
Julia set.  The available
quantitative local
straightening formulation
\cite[Lemma 4.33]{LyubichNotes}
depends on its chosen
modulus, but furnishes
no explicit numerical
global marked slit
dilatation.  Neither
holomorphic-motion
extension
\cite{BersRoyden1986,Slodkowski1991}
nor Beltrami pullback
manufactures a missing
external prime-end marking
or an identification of
\emph{every} KKL A.1
and A.3.2 Wave class.

Here is an exact
conditional implication,
which makes the missing
hypothesis testable.
Let $\mathcal A_j$ be
the entire admissible
parent Wave class
on a physical depth-$j$
marked slit and
$\mathcal A_{\rm ref}$
a specified single-level
reference class.  Assume
for some \emph{common}
$C\ge1$:
\begin{itemize}
\item for every $\Gamma\in\mathcal A_j$
there is $\gamma\in\mathcal A_{\rm ref}$
with $\W(\Gamma)\le C\W(\gamma)$;
\item for every $\gamma\in\mathcal A_{\rm ref}$
there is $\Gamma\in\mathcal A_j$
with $\W(\gamma)\le C\W(\Gamma)$.
\end{itemize}
Taking suprema \emph{separately}
proves
\begin{equation}\label{eq:depth}
 \frac{W_{\rm ref}^*}{C}
 \le W_j^*\le C W_{\rm ref}^*.
\end{equation}
This \emph{two-sided
cofinal domination}
is strictly weaker than
the existence of a
global QC map.  A
globally defined real
equivariant $K'$-QC
marked conjugacy that
extends to the \emph{entire}
physical prime-end
completions, respects
$G_*$ and the set-valued
$\iota^*$, and
bijects all A.1/A.3.2
Wave classes would
imply it with $C=K'$
by extremal-length
quasi-invariance
\cite{AhlforsQC1966}.
\emph{Such a map has
not been constructed.}
If one \emph{assumes}
it with a depth-independent
$K'$, a general
single-level reference
gives
\[
 0.25/K'<W_j^*<0.78K',
\]
and a separately certified
reference with period
$p_{\rm ref}\ge13$ gives
$0.26/K'<W_j^*<0.78K'$.
These are \emph{conditional}
at every depth, with no
$K'^j$ factor because
the comparison is directly
to the reference.
For a positive \emph{matched}
parent and collared
four-push descendant,
their ratio is distorted
by at most $K'^2$
in either direction;
this does not apply
to two unrelated
optimizing suprema.
At infinite depth the
analogous liminf/limsup
bounds are \emph{closed},
and no limiting marked
slit or Wave convergence
is constructed.

A concrete obstruction
to \emph{deducing}
global marked $K'$
from local straightening
is furnished by
\[
 A_1=\{1<|z|<e^{2\pi}\},\qquad
 A_N=\{1<|z|<e^{2\pi N}\}.
\]
Both have an identical
real-symmetric conformal
collar near the inner
boundary, but their
full marked annuli
have moduli $1$ and
$N$.  Any QC map of
\emph{both marked
boundary components}
has $K\ge N$ by
modulus distortion.
This is a countermodel
to the \emph{logical
inference from local
data}, not a
counterexample in the
actual primitive-tip
KKL family.
The all-depth global
marked-QC transport
and an effective
$K'(\mu)$ are open.

\section{Limitations and disclosure}
The proved mathematical
statements are about
the \emph{selected
real} quadratic family
and the exact physical
slit surfaces.  The
local parameter
uniqueness does not
extend to all real
solutions or all
parameter components.
Neither the positive
unbuffered independent
Wave defect nor the
finite edge matrix
gives a positive
buffered WAD coefficient
or KKL \S5--6
amplification.  The
four-boundary
\emph{inner} buffered
support, maximizing
Waves, convergence of
canonical unrestricted
widths, the
four-push crest limit,
an all-word or
infinite-depth optimal
Wave theorem, and
uniform global
marked quasiconformal
transport remain open.
The family is real
and MLC in this
setting was already
settled by classical
theory; \emph{no
new MLC result} is
proved here.
The computation is
not a replacement
for independent
review of the
cited analytic proofs
or of this manuscript.
This text was
drafted and assembled
using AI assistance
from merged public
audit sources; it is
not yet independently
reviewed or
submitted as a
peer-reviewed
publication.
